%% papericai.tex — Springer LNCS Proceedings Paper
%% Auto-generated from Master's Thesis
%% Target: 12–16 pages
%%
%% Thesis: "Accent Based Model for Speech Spoofing Countermeasures"
%% Author: Jaime Arturo Hurtado Romero
%% Advisor: Prof. Ruben Francisco Manrique
%% FLAG Lab, Universidad de los Andes, Bogotá, Colombia
%%
%% Compilation note: SVG figures require Inkscape to be installed and
%%   the svg LaTeX package (included below). If Inkscape is unavailable,
%%   replace \includesvg{...} with \includegraphics{...} and export SVGs
%%   to PDF/PNG externally.

\documentclass{llncs}

\usepackage{graphicx}
\usepackage{amsmath}
\usepackage{amssymb}
\usepackage{multirow}
\usepackage{xcolor}
\usepackage{array}  
\usepackage{svg}
\svgpath{{./svg/}}
\usepackage{url}

\begin{document}

\title{Accent-Based Evaluation of Speech Anti-Spoofing 
       Countermeasures Across Multiple Languages}

\author{Jaime Arturo Hurtado Romero\inst{1} \and Tomas Acosta \and 
        Ruben Francisco Manrique\inst{1}}

\institute{FLAG Lab, Department of Systems and Computing Engineering,\\
           Universidad de los Andes, Bogot\'{a}, Colombia\\
           \email{\{ja.hurtado, t.acosta, rf.manrique\}@uniandes.edu.co}}

\maketitle

%% ----------------------------------------------------------------
\begin{abstract}
Speech-based biometric systems face increasing threats from sophisticated
spoofing attacks enabled by advances in deep generative modelling.  A
critical yet under-studied limitation is the severe performance degradation
that occurs when anti-spoofing countermeasure models trained on English
corpora are applied to Spanish speech, owing to systematic accent-induced
differences in prosodic and spectral features.  This paper presents a
systematic evaluation of 48 anti-spoofing architectures formed by crossing
four front-end feature extractors---Mel-frequency Cepstral Coefficients
(MFCC), Linear Frequency Cepstral Coefficients (LFCC), Constant-Q Cepstral
Coefficients (CQCC), and Spectrogram---with three Light Convolutional Neural
Network (LCNN) back-ends (Attention, Trim-pad, LSTM-sum) and four loss
functions.  Models are trained and evaluated on the ASVspoof~2019 English
database and the HABLA Spanish multi-accent database, as well as a combined
bilingual corpus.  Each configuration is run six times with different random
seeds, and results are reported as mean Equal Error Rate (EER).  Joint
English--Spanish training achieves near-zero EER on Spanish (best:
MFCC-LCNN-LSTM-sum with AM-Softmax at $0.09\%$) while sacrificing fewer
than one percentage point on English, demonstrating that a single bilingualy
trained model can generalise across languages.  An accent-granularity
analysis over five Latin-American accents further reveals that the optimal
architecture varies per accent, motivating accent-adaptive or ensemble
strategies in future work.
\end{abstract}

\keywords{speech anti-spoofing \and accent diversity \and cross-language evaluation
          \and MFCC \and LFCC \and LCNN \and ASVspoof \and HABLA}

%% ----------------------------------------------------------------
\section{Introduction}

Biometric systems that rely on speech for identity verification are exposed
to a growing range of spoofing attacks.  \textit{Speech synthesis}
(text-to-speech, TTS) produces spoken output from written input, while
\textit{voice conversion} (VC) modifies non-linguistic characteristics of a
waveform to impersonate a target speaker~\cite{stylianou1998continuous}.
Both attack types can be mounted at the logical access point of an automatic
speaker verification (ASV) system and may lead to unauthorised access,
financial fraud, or identity theft~\cite{xue2023physiological}.  The
proliferation of Generative Adversarial Networks (GANs)~\cite{goodfellow2020generative}
and Variational Autoencoders (VAEs)~\cite{kingma2013auto} has dramatically
lowered the barrier to generating perceptually convincing synthetic speech,
amplifying the urgency of robust anti-spoofing countermeasures
(CMs)~\cite{kaur2014review}.

The \textit{ASVspoof} challenge series~\cite{yamagishi2019asvspoof,
delgado2021asvspoof} has driven most recent progress in anti-spoofing.
However, all ASVspoof benchmark datasets are exclusively in English.
Prior work has shown that state-of-the-art CMs trained on English data
degrade severely when evaluated on Spanish speech~\cite{tamayo2022voice},
with Equal Error Rates exceeding $30\%$ even for strong baselines.  This
degradation is linked to accent-driven differences in vowel space, formant
structure, and speaking rate that directly affect the cepstral features
commonly employed as front-ends.  Despite this, there is limited research
that \textit{systematically} quantifies how specific front-end/back-end
combinations respond to language and accent variation.

\textbf{Contributions.} This work makes the following contributions:

\begin{enumerate}
  \item A systematic comparison of four front-end feature extractors combined
        with three LCNN back-end variants and four loss functions (48
        architecture combinations), each trained and evaluated across three
        dataset conditions.
  \item Cross-language experiments evaluating models trained on
        English, Spanish, and a combined bilingual corpus on both languages.
  \item An accent-granularity analysis over five Spanish regional accents
        (Argentina, Chile, Colombia, Peru, Venezuela), revealing
        accent-specific performance trends that cannot be inferred from
        aggregate metrics.
  \item A filter-bank scale analysis demonstrating that descendant
        (low-frequency emphasis) scaling consistently benefits cepstral
        front-ends, consistent with the importance of low-frequency voice
        formants.
\end{enumerate}

%% ----------------------------------------------------------------
\section{Related Work}

\subsection{Anti-Spoofing Datasets}

The ASVspoof benchmarks are the primary evaluation corpora for logical-access
(LA) and physical-access (PA) spoofing.  ASVspoof~2019~\cite{yamagishi2019asvspoof}
covers LA and PA attacks, employing 19 TTS/VC systems in training and 48 in
evaluation; ASVspoof~2021~\cite{delgado2021asvspoof} further extends to
deepfake (DF) speech.  All ASVspoof datasets are English only.  The
\textit{HABLA Spoof} corpus~\cite{tamayo2022voice} addresses the Spanish
gap: it contains synthesised and converted speech from speakers representing
five Latin-American accents, enabling accent-granularity evaluation of
anti-spoofing systems.  Table~\ref{table:datasets-taxonomy} summarises the
main datasets considered in this study.

\begin{table}[htbp]
\centering
\caption{Anti-spoofing dataset taxonomy. LA = Logical Access, PA = Physical
         Access, DF = Deepfake; EN = English, ES = Spanish.}
\label{table:datasets-taxonomy}
\setlength\extrarowheight{-2pt}
\begin{tabular}{lcl}
\hline
\textbf{Dataset}                           & \textbf{Lang.} & \textbf{Attacks} \\
\hline
HABLA Spoof~\cite{tamayo2022voice}         & ES & LA, DF       \\
ASVspoof 2019~\cite{yamagishi2019asvspoof} & EN & LA, PA       \\
ASVspoof 2021~\cite{delgado2021asvspoof}   & EN & LA, PA, DF   \\
SAS~\cite{wu2016anti}                      & EN & LA           \\
FoR~\cite{reimao2019dataset}               & EN & LA           \\
VCTK~\cite{yamagishi2019cstr}              & EN & LA           \\
\hline
\end{tabular}
\end{table}

\subsection{Front-End Feature Extractors}

Digital Signal Processing (DSP)-based front-ends extract a compact acoustic
feature sequence from each utterance.  MFCC~\cite{davis1980comparison}
applies a Mel-scale filter bank followed by a Discrete Cosine Transform
(DCT), producing perceptually motivated coefficients compact in the low
frequencies.  LFCC~\cite{zhou2011linear} replaces the Mel scale with a
uniform linear-scale filter bank, treating all frequencies equally.
CQCC pairs the Constant-Q Transform with cepstral analysis, providing
geometrically spaced frequency resolution suited to capturing spectral
artefacts in synthesised speech.  Spectrogram features retain the full
Short-Time Fourier Transform (STFT) magnitude, serving as the richest
but highest-dimensional representation.  Both LFCC and CQCC were the
official ASVspoof baselines~\cite{wang2022practical}.

\subsection{Back-End Classifiers}

The back-end maps the front-end feature sequence to a spoofing score.
Light Convolutional Neural Networks (LCNNs)~\cite{lavrentyeva2017audio}
use Max-Feature-Map (MFM) activations to suppress irrelevant feature
channels and have become a strong standard baseline.
LSTM-augmented LCNNs~\cite{wang2021comparative} further model long-range
temporal dependencies.  End-to-end models such as RawNet2~\cite{tak2021end}
operate directly on raw waveforms, while classical Gaussian Mixture Models
(GMMs) remain competitive for PA scenarios~\cite{cai2017countermeasures}.
Table~\ref{table:anti-spoofing-taxonomy} lists representative
front-end/back-end pairings from the literature.

\begin{table}[htbp]
\centering
\caption{Representative anti-spoofing model taxonomy from the literature.}
\label{table:anti-spoofing-taxonomy}
\setlength\extrarowheight{-2pt}
\begin{tabular}{lll}
\hline
\textbf{Front-end}         & \textbf{Back-end}        & \textbf{Reference}          \\
\hline
Spectrogram                & CNN + RNN                & \cite{zhang2017investigation} \\
LFCC / IMFCC               & GMM \& TDNN              & \cite{wu2014study}            \\
CQCC + Spectrogram         & GMM \& ResNet            & \cite{cai2017countermeasures} \\
Spectrogram + LCNN         & RNN                      & \cite{lavrentyeva2017audio}   \\
Spectrogram                & SENet34                  & \cite{lai2019assert}          \\
Raw wave                   & RawNet2                  & \cite{tak2021end}             \\
\hline
\end{tabular}
\end{table}

\subsection{Research Gap}

Existing anti-spoofing literature is concentrated on English corpora and
single-language settings.  Tamayo~et~al.~\cite{tamayo2022voice} demonstrated
the cross-lingual drop but did not systematically evaluate front-end/back-end
combinations or filter-bank scales.  The effect of regional accents on
feature extraction and the degree to which multilingual training mitigates
language mismatch remain largely open questions.  Our work directly addresses
this gap.

%% ----------------------------------------------------------------
\section{Methodology}

\subsection{Experimental Framework}

Each anti-spoofing model is composed of a \textit{front-end} that extracts
acoustic features and a \textit{back-end} that classifies the utterance as
bona fide or spoofed.  We systematically combine all front-end and back-end
variants in Table~\ref{table:antispoofing-modules}, yielding $4 \times 3 \times
4 = 48$ configurations per training condition.

\begin{table}[htbp]
\centering
\caption{Anti-spoofing architecture modules explored in this study.}
\label{table:antispoofing-modules}
\setlength\extrarowheight{-2pt}
\begin{tabular}{lll}
\hline
\textbf{Front-End} & \textbf{Back-End}    & \textbf{Loss / Activation} \\
\hline
CQCC               & LCNN-Attention       & MSE for P2S-Grad \\
MFCC               & LCNN-trim-pad        & AM-Softmax       \\
LFCC               & LCNN-LSTM-sum        & OC-Softmax       \\
Spectrogram        &  ---                 & Sigmoid          \\
\hline
\end{tabular}
\end{table}

\subsection{Front-End Feature Extraction}

\textbf{MFCC and LFCC} are cepstral coefficient representations computed
frame-by-frame via Discrete Fourier Transform (DFT), log-magnitude, and DCT.
MFCC uses a Mel-scale (logarithmically spaced) filter bank that mimics human
auditory sensitivity; LFCC uses a linearly spaced bank that weights all
frequencies equally~\cite{zhou2011linear}.  We vary the filter-bank
\textit{scale response} in three modes that modulate how gain changes across
frequency: \textit{ascendant} (gain increases toward high frequencies),
\textit{constant} (flat gain), and \textit{descendant} (gain decreases toward
high frequencies, emphasising low-frequency formants).

\textbf{CQCC} applies the Constant-Q Transform before cepstral analysis,
providing finer frequency resolution at low frequencies where anti-spoofing
artefacts are often present.

\textbf{Spectrogram} retains the full STFT magnitude as a 2-D image fed
directly into the LCNN.

\subsection{LCNN Back-End Architecture}

All back-ends share a common LCNN feature-extraction stem followed by one of
three utterance-level pooling strategies.  \textit{Attention-based} pooling
learns to weight frames by relevance; \textit{trim-pad} forces fixed-length
input; \textit{LSTM-sum} models temporal dynamics with a bidirectional LSTM
before summing.  The pooled embedding is fed to a classification head whose
output distribution is shaped by the chosen loss function: AM-Softmax,
OC-Softmax, MSE-based P2S-Grad, or Sigmoid cross-entropy.

The LCNN stem architecture is detailed in Table~\ref{table:lcnn-architecture}.
It employs MFM activations after every convolution, batch normalisation where
indicated, and max-pooling for spatial downsampling.

\begin{table}[htbp]
\centering
\caption{LCNN stem architecture shared across all back-end variants.
         MFM = Max-Feature-Map, BN = Batch Normalisation.}
\label{table:lcnn-architecture}
\setlength\extrarowheight{-4pt}
\footnotesize
\begin{tabular}{lcc}
\hline
\textbf{Layer}         & \textbf{Filter/Stride}      & \textbf{Output}          \\
\hline
Conv + MFM             & $5{\times}5$ / $1{\times}1$ & $864{\times}600{\times}32$ \\
MaxPool                & $2{\times}2$ / $2{\times}2$ & $432{\times}300{\times}32$ \\
\hline
Conv + MFM + BN        & $1{\times}1$ / $1{\times}1$ & $432{\times}300{\times}32$ \\
Conv + MFM             & $3{\times}3$ / $1{\times}1$ & $432{\times}300{\times}48$ \\
MaxPool + BN           & $2{\times}2$ / $2{\times}2$ & $216{\times}150{\times}48$ \\
\hline
Conv + MFM + BN        & $1{\times}1$ / $1{\times}1$ & $216{\times}150{\times}48$ \\
Conv + MFM             & $3{\times}3$ / $1{\times}1$ & $216{\times}150{\times}64$ \\
MaxPool                & $2{\times}2$ / $2{\times}2$ & $108{\times}75{\times}64$  \\
\hline
4$\times$(Conv+MFM+BN) & $1{\times}1$--$3{\times}3$ / $1{\times}1$ & $108{\times}75{\times}32$ \\
MaxPool                & $2{\times}2$ / $2{\times}2$ & $54{\times}38{\times}32$   \\
\hline
\end{tabular}
\end{table}

\subsection{Datasets and Partitions}

Two anti-spoofing corpora are used; partition sizes are listed in
Table~\ref{table:datasets-partitions}.

\textbf{English (ASVspoof~2019 LA)}~\cite{yamagishi2019asvspoof}: Standard
logical-access split with 19 TTS/VC systems in training data and 48 in
the evaluation set.

\textbf{Spanish (HABLA Spoof)}~\cite{tamayo2022voice}: Latin-American
multi-accent corpus covering five countries (Argentina, Chile, Colombia,
Peru, Venezuela).  Speakers are assigned exclusively to a single partition
and are never shared across training, validation, and test splits.

\textbf{Combined}: The English and Spanish training sets are concatenated
to form a 60\,064-utterance bilingual training corpus; evaluation is
performed separately on each language's test partition.

\begin{table}[htbp]
\centering
\caption{Dataset partition sizes (number of utterances).}
\label{table:datasets-partitions}
\setlength\extrarowheight{-2pt}
\begin{tabular}{lccc}
\hline
\textbf{Training set} & \textbf{Train} & \textbf{Dev} & \textbf{Test} \\
\hline
English               & 25\,380 & 24\,840 & 71\,237 \\
Spanish               & 34\,684 &  3\,658 &  4\,809 \\
Combined              & 60\,064 & 28\,498 & --- \\
\hline
\end{tabular}
\end{table}

\subsection{Training Protocol}

All models are trained for up to 100 epochs with an Adam
optimiser~\cite{wang2021comparative} ($\beta_1=0.9$, $\beta_2=0.999$,
$\epsilon=10^{-8}$).  The initial learning rate of $3\times10^{-4}$ is
halved every ten epochs.  Mini-batches of 64 utterances are sampled by
similar duration.  Early stopping terminates training after 50 epochs
without improvement on the development loss; the checkpoint with the
lowest development loss is retained.  No voice activity detection or
feature normalisation is applied.  To account for variance due to random
initialisation, every configuration is trained \textbf{six} times with
different random seeds; results are reported as \emph{mean $\pm$ standard
deviation EER~(\%)}.

Figure~\ref{fig:loss_curve} shows a representative training curve,
illustrating the interplay between the learning-rate schedule, the
best-model callback, and early stopping.

\begin{figure}[htbp]
    \centering
    \includegraphics[width=0.85\textwidth]{figures/loss_curve}
    \caption{Training curve for the LFCC-LCNN-Attention-AM model (logarithmic
             scale).  The best-model callback retains the minimum dev-loss
             checkpoint; early stopping triggers at epoch~72, 50 epochs after
             the last improvement.}
    \label{fig:loss_curve}
\end{figure}

\subsection{Evaluation Metric}

Models are evaluated using the \textit{Equal Error Rate} (EER), the
operating point at which the False Rejection Rate (FRR) equals the False
Acceptance Rate (FAR):

\begin{equation}
    \text{EER} = \frac{\text{FRR}(\tau') + \text{FAR}(\tau')}{2},
    \qquad
    \tau' = \underset{\tau}{\arg\min}\,|\text{FRR}(\tau) - \text{FAR}(\tau)|.
    \label{eq:eer}
\end{equation}

A lower EER indicates better discrimination between bona-fide and
spoofed utterances.

%% ----------------------------------------------------------------
\section{Experiments and Results}

We structure the experiments around five conditions, systematically varying
the training language and the evaluation language.  The code and trained
model weights are publicly available~\cite{MisisThesisCode}.

\subsection{English Training, English Evaluation}

Table~\ref{table:mean-train-en-eval-en} and
Figure~\ref{fig:heatmap-train-en-eval-en} report results when all models
are trained and evaluated on the English partitions.  Performance is
satisfactory across all 48 configurations.  Cepstral front-ends (LFCC and
MFCC) consistently outperform CQCC and Spectrogram.

The overall best configuration is \textbf{LFCC-LCNN-LSTM-sum with P2S-Grad}
at $2.67\%$ EER ($\sigma=0.60$).  P2S-Grad loss dominates across front-ends;
OC-Softmax is the second strongest.  Spectrogram front-ends show higher
variance than cepstral ones, and CQCC performs comparably to Spectrogram in
this English-only setting.

\begin{figure}[htbp]
    \centering
    \includesvg[inkscapelatex=false, width=\textwidth]{english_heatmap}
    \caption{EER (\%) heat map for models trained and evaluated on the English
             partition.  Darker cells indicate higher EER.  Rows are sorted
             in ascending EER order.}
    \label{fig:heatmap-train-en-eval-en}
\end{figure}

\begin{table}[htbp]
\centering
\caption{Mean and std of EER (\%) for English train / English eval.
         Best per back-end in \textbf{bold}; overall best in
         \textcolor{red}{\textbf{red}}.}
\label{table:mean-train-en-eval-en}
\setlength\extrarowheight{-8pt}
\footnotesize
\begin{tabular}{llcccccc}
\hline
 & & \multicolumn{2}{c}{Attention} & \multicolumn{2}{c}{Trim-pad}
   & \multicolumn{2}{c}{LSTM-sum} \\
\textbf{Front} & \textbf{Activation} &
  \textbf{Mn} & \textbf{Sd} & \textbf{Mn} & \textbf{Sd} &
  \textbf{Mn} & \textbf{Sd} \\
\hline
\textbf{CQCC}
  & AM-softmax & 7.58 & 0.69 & 7.49 & 0.73 & 7.81 & 0.73 \\
  & P2SGrad    & 6.65 & 0.33 & \textbf{6.00} & \textbf{0.19} & 6.20 & 0.14 \\
  & OC-softmax & \textbf{6.40} & \textbf{0.30} & 6.12 & 0.12 & 6.63 & 0.32 \\
  & Sigmoid    & 8.08 & 1.01 & 6.23 & 0.32 & \textbf{6.14} & \textbf{0.22} \\
\hline
\textbf{Spec.}
  & AM-softmax & 5.35 & 1.53 & 5.52 & 0.56 & 4.67 & 0.76 \\
  & P2SGrad    & 5.68 & 0.71 & \textbf{3.79} & \textbf{0.71} & \textbf{3.48} & \textbf{0.82} \\
  & OC-softmax & 5.67 & 1.59 & 5.10 & 0.70 & 4.13 & 0.83 \\
  & Sigmoid    & \textbf{4.81} & \textbf{0.98} & 4.50 & 0.96 & 4.19 & 0.79 \\
\hline
\textbf{MFCC}
  & AM-softmax & 4.46 & 0.37 & 6.97 & 1.95 & 4.73 & 2.30 \\
  & P2SGrad    & 3.04 & 0.51 & \textbf{2.88} & \textbf{0.34} & \textbf{2.83} & \textbf{0.29} \\
  & OC-softmax & 3.71 & 1.37 & 4.49 & 1.09 & 3.51 & 0.93 \\
  & Sigmoid    & \textbf{3.50} & \textbf{0.17} & 3.52 & 0.61 & 3.47 & 0.94 \\
\hline
\textbf{LFCC}
  & AM-softmax & 4.00 & 0.33 & 5.58 & 0.61 & 3.25 & 1.11 \\
  & P2SGrad    & 3.48 & 0.78 & 2.88 & 0.41 & \textcolor{red}{\textbf{2.67}} & \textcolor{red}{\textbf{0.60}} \\
  & OC-softmax & \textbf{2.96} & \textbf{0.35} & \textbf{2.88} & \textbf{0.23} & 2.80 & 0.61 \\
  & Sigmoid    & 3.73 & 0.35 & 3.04 & 0.36 & 3.33 & 0.31 \\
\hline
\end{tabular}
\end{table}

\subsection{Cross-Language Evaluation: English Models on Spanish}

Applying the English-trained models directly to the Spanish test partition
reveals severe degradation: \emph{all} 48 configurations exceed $30\%$ EER,
and most exceed $40\%$.  Among front-ends, Spectrogram degrades the least
(best mean $\approx 36\%$), while CQCC degrades the most ($43\text{--}48\%$)
despite being the strongest English front-end.  This inversion suggests that
CQCC captures English-specific phonetic patterns that do not transfer to
Spanish prosody.  MFCC and LFCC, despite being strong English baselines,
yield cross-lingual EERs of $34\text{--}55\%$, confirming that language-specific
training is essential for robust performance.

\subsection{Spanish Training, Spanish Evaluation}

Training and evaluating exclusively on Spanish data yields excellent results
for most architectures (Table~\ref{table:mean-train-es-eval-es}).  The best
configuration is \textbf{LFCC-LCNN-Attention with P2S-Grad} at $0.14\%$ EER
($\sigma=0.09$).  Spectrogram front-ends are the notable exception,
consistently yielding EER of $12\text{--}20\%$, suggesting that synthetic
Spanish speech lacks the texture-level discriminability that Spectrogram
features exploit in English.

Comparing Table~\ref{table:mean-train-es-eval-es} with the English baseline,
the \emph{architecture ranking changes}: LFCC-LCNN-Attention/P2S-Grad is
best for Spanish whereas LFCC-LCNN-LSTM-sum/P2S-Grad was best for English.
This language-specific sensitivity motivates the combined-training approach.

\begin{table}[htbp]
\centering
\caption{Mean and std of EER (\%) for Spanish train / Spanish eval (best
         activation per front-end/back-end shown).  Overall best in
         \textcolor{red}{\textbf{red}}.}
\label{table:mean-train-es-eval-es}
\setlength\extrarowheight{-8pt}
\footnotesize
\begin{tabular}{llcccccc}
\hline
 & & \multicolumn{2}{c}{Attention} & \multicolumn{2}{c}{Trim-pad}
   & \multicolumn{2}{c}{LSTM-sum} \\
\textbf{Front} & \textbf{Activation} &
  \textbf{Mn} & \textbf{Sd} & \textbf{Mn} & \textbf{Sd} &
  \textbf{Mn} & \textbf{Sd} \\
\hline
\textbf{CQCC}
  & AM-softmax & 1.35 & 1.16 & 4.15 & 2.15 & 2.54 & 1.34 \\
  & P2SGrad    & \textbf{0.61} & \textbf{0.18} & 3.87 & 1.82 & 2.69 & 1.27 \\
  & OC-softmax & 2.69 & 2.15 & 4.73 & 1.87 & \textbf{1.29} & \textbf{0.48} \\
  & Sigmoid    & 1.29 & 0.63 & \textbf{3.36} & \textbf{1.93} & 2.53 & 1.43 \\
\hline
\textbf{Spec.}
  & AM-softmax & 17.15 & 3.33 & 17.98 & 4.05 & 17.30 & 1.50 \\
  & P2SGrad    & \textbf{12.14} & \textbf{3.02} & \textbf{17.24} & \textbf{1.43} & 19.50 & 4.96 \\
  & OC-softmax & 15.83 & 2.50 & 17.41 & 1.83 & \textbf{15.76} & \textbf{1.90} \\
  & Sigmoid    & 20.33 & 3.61 & 20.23 & 3.50 & 19.58 & 2.55 \\
\hline
\textbf{MFCC}
  & AM-softmax & 1.10 & 1.28 & 3.79 & 3.11 & 0.65 & 0.57 \\
  & P2SGrad    & \textbf{0.44} & \textbf{0.77} & \textbf{3.10} & \textbf{1.84} & 1.21 & 1.20 \\
  & OC-softmax & 0.77 & 0.73 & 3.42 & 4.82 & 1.38 & 1.14 \\
  & Sigmoid    & 1.74 & 2.34 & 3.25 & 2.83 & \textbf{0.42} & \textbf{0.31} \\
\hline
\textbf{LFCC}
  & AM-softmax & 1.40 & 0.90 & 3.58 & 3.05 & \textbf{0.45} & \textbf{0.27} \\
  & P2SGrad    & \textcolor{red}{\textbf{0.14}} & \textcolor{red}{\textbf{0.09}} & 3.55 & 2.80 & 2.32 & 1.57 \\
  & OC-softmax & 0.68 & 0.90 & 3.89 & 4.70 & 1.30 & 1.71 \\
  & Sigmoid    & 0.49 & 0.50 & 4.07 & 5.36 & 0.57 & 0.55 \\
\hline
\end{tabular}
\end{table}

\subsection{Combined Bilingual Training}

Training on the merged English--Spanish corpus provides the most practically
robust models.  Figure~\ref{fig:heatmap-train-un-eval-es} and
Table~\ref{table:mean-train-un-eval-es} show results on the Spanish
evaluation partition; Table~\ref{table:mean-train-un-eval-en} shows results
on the English partition.

\begin{figure}[htbp]
    \centering
    \includesvg[inkscapelatex=false, width=\textwidth]{united-eval-lat_heatmap}
    \caption{EER (\%) heat map for models trained on the combined
             English--Spanish corpus and evaluated on Spanish.  Darker cells
             indicate higher EER.}
    \label{fig:heatmap-train-un-eval-es}
\end{figure}

MFCC and LFCC front-ends achieve near-zero EER ($\leq 0.5\%$) on Spanish
evaluation, with the top result being \textbf{MFCC-LCNN-LSTM-sum with
AM-Softmax} at $\mathbf{0.09\%}$ ($\sigma=0.04$).  CQCC also performs well
($\leq 1.3\%$).  Spectrogram remains the weakest front-end, mirroring the
Spanish-only experiment.

On the English evaluation partition (Table~\ref{table:mean-train-un-eval-en}),
performance decreases by only $\sim$1 percentage point relative to the
English-only baseline.  The best bilingual English configuration is
MFCC-LCNN-trim-pad with P2S-Grad at $3.64\%$.  This near-parity on English
together with near-zero EER on Spanish confirms that joint training achieves
genuine bilingual generalisation.

\begin{table}[htbp]
\centering
\caption{Mean and std of EER (\%) for combined-trained models, Spanish eval.
         Best per back-end in \textbf{bold}; overall best in
         \textcolor{red}{\textbf{red}}.}
\label{table:mean-train-un-eval-es}
\setlength\extrarowheight{-8pt}
\footnotesize
\begin{tabular}{llcccccc}
\hline
 & & \multicolumn{2}{c}{Attention} & \multicolumn{2}{c}{Trim-pad}
   & \multicolumn{2}{c}{LSTM-sum} \\
\textbf{Front} & \textbf{Activation} &
  \textbf{Mn} & \textbf{Sd} & \textbf{Mn} & \textbf{Sd} &
  \textbf{Mn} & \textbf{Sd} \\
\hline
\textbf{CQCC}
  & AM-softmax & 1.16 & 0.43 & 1.15 & 0.38 & 0.97 & 0.24 \\
  & P2SGrad    & 0.96 & 0.36 & \textbf{0.89} & \textbf{0.25} & \textbf{0.49} & \textbf{0.18} \\
  & OC-softmax & \textbf{0.86} & \textbf{0.48} & 0.97 & 0.45 & 1.25 & 1.40 \\
  & Sigmoid    & 2.98 & 1.29 & 1.35 & 0.25 & 0.82 & 0.47 \\
\hline
\textbf{Spec.}
  & AM-softmax & 17.10 & 1.50 & 20.82 & 2.72 & 15.22 & 1.44 \\
  & P2SGrad    & \textbf{16.37} & \textbf{2.66} & 16.46 & 3.61 & 15.35 & 2.87 \\
  & OC-softmax & 17.88 & 7.00 & \textbf{17.32} & \textbf{1.62} & \textbf{14.92} & \textbf{1.83} \\
  & Sigmoid    & 16.65 & 2.40 & 18.58 & 1.35 & 16.96 & 2.74 \\
\hline
\textbf{MFCC}
  & AM-softmax & 0.31 & 0.22 & 0.42 & 0.17 & \textcolor{red}{\textbf{0.09}} & \textcolor{red}{\textbf{0.04}} \\
  & P2SGrad    & \textbf{0.17} & \textbf{0.09} & \textbf{0.17} & \textbf{0.09} & 0.35 & 0.33 \\
  & OC-softmax & 0.40 & 0.34 & 0.29 & 0.03 & 0.45 & 0.47 \\
  & Sigmoid    & 0.55 & 0.36 & 0.38 & 0.13 & 0.45 & 0.49 \\
\hline
\textbf{LFCC}
  & AM-softmax & 0.25 & 0.14 & \textbf{0.30} & \textbf{0.09} & 0.37 & 0.40 \\
  & P2SGrad    & \textbf{0.19} & \textbf{0.10} & 0.30 & 0.15 & 0.30 & 0.26 \\
  & OC-softmax & 0.42 & 0.36 & 0.34 & 0.13 & 0.24 & 0.12 \\
  & Sigmoid    & 0.95 & 0.87 & 0.28 & 0.11 & \textbf{0.18} & \textbf{0.14} \\
\hline
\end{tabular}
\end{table}

\begin{table}[htbp]
\centering
\caption{Mean and std of EER (\%) for combined-trained models, English eval.
         Best per back-end in \textbf{bold}; overall best in
         \textcolor{red}{\textbf{red}}.}
\label{table:mean-train-un-eval-en}
\setlength\extrarowheight{-8pt}
\footnotesize
\begin{tabular}{llcccccc}
\hline
 & & \multicolumn{2}{c}{Attention} & \multicolumn{2}{c}{Trim-pad}
   & \multicolumn{2}{c}{LSTM-sum} \\
\textbf{Front} & \textbf{Activation} &
  \textbf{Mn} & \textbf{Sd} & \textbf{Mn} & \textbf{Sd} &
  \textbf{Mn} & \textbf{Sd} \\
\hline
\textbf{CQCC}
  & AM-softmax & 6.93 & 2.54 & 6.43 & 0.40 & 6.65 & 0.71 \\
  & P2SGrad    & 5.77 & 0.33 & \textbf{5.77} & \textbf{0.11} & \textbf{5.71} & \textbf{0.31} \\
  & OC-softmax & 6.41 & 0.36 & 5.83 & 0.09 & 5.77 & 0.27 \\
  & Sigmoid    & 9.10 & 0.75 & 5.87 & 0.22 & \textbf{5.51} & \textbf{0.21} \\
\hline
\textbf{Spec.}
  & AM-softmax & 11.19 & 1.35 & 9.95 & 0.51 & 10.14 & 0.54 \\
  & P2SGrad    & 10.45 & 1.57 & \textbf{9.16} & \textbf{0.42} & \textbf{9.06} & \textbf{0.49} \\
  & OC-softmax & 10.95 & 1.23 & 10.13 & 0.50 & 9.98 & 1.23 \\
  & Sigmoid    & \textbf{10.26} & \textbf{0.74} & 10.29 & 0.72 & 10.51 & 0.99 \\
\hline
\textbf{MFCC}
  & AM-softmax & 4.96 & 1.27 & 4.43 & 0.67 & 4.62 & 0.71 \\
  & P2SGrad    & 4.07 & 0.47 & \textcolor{red}{\textbf{3.64}} & \textcolor{red}{\textbf{0.42}} & 4.58 & 1.02 \\
  & OC-softmax & \textbf{3.86} & \textbf{0.66} & 3.75 & 0.52 & 4.97 & 1.18 \\
  & Sigmoid    & 4.64 & 0.77 & 3.66 & 0.37 & \textbf{4.26} & \textbf{0.84} \\
\hline
\textbf{LFCC}
  & AM-softmax & 4.67 & 0.73 & 4.54 & 0.64 & 4.43 & 0.79 \\
  & P2SGrad    & \textbf{4.37} & \textbf{0.67} & 3.98 & 0.17 & \textbf{4.28} & \textbf{0.42} \\
  & OC-softmax & 5.00 & 1.31 & \textbf{3.78} & \textbf{0.44} & 4.50 & 1.22 \\
  & Sigmoid    & 5.11 & 0.74 & 3.79 & 0.49 & 4.24 & 0.51 \\
\hline
\end{tabular}
\end{table}

\subsection{Accent Analysis}

Using the bilingualy trained models, we break down Spanish evaluation EER
by Latin-American accent.  Figure~\ref{fig:accent-MFCC} shows results for
the MFCC front-end; LFCC yields qualitatively similar patterns.

\begin{figure}[htbp]
    \centering
    \includesvg[inkscapelatex=false, width=\textwidth]{accent_MFCC}
    \caption{Accent-level EER (\%) for MFCC front-end models trained on the
             combined corpus, disaggregated by Latin-American accent.}
    \label{fig:accent-MFCC}
\end{figure}

Key observations:

\begin{itemize}
  \item \textbf{Chile and Peru} achieve consistently near-zero EER across
        architectures, indicating that these accents are well-modelled by
        general cepstral features.
  \item \textbf{Argentina} performs well but benefits disproportionately from
        LCNN-Attention with P2S-Grad, suggesting that its prosodic signature
        aligns with the attention mechanism's learned weighting.
  \item \textbf{Colombia} yields higher EER overall, attributable to the
        notable intra-country accent diversity that is not fully captured by
        the HABLA corpus.
  \item \textbf{Venezuela} is the most challenging accent: LCNN-LSTM-sum
        outperforms LCNN-Attention here, the \emph{reverse} of the pattern
        for other accents.  This suggests temporally distributed cues specific
        to Venezuelan prosody that LSTM pooling is better suited to capture.
\end{itemize}

These findings confirm that no single architecture uniformly dominates across
all accents, motivating accent-adaptive or multi-task learning approaches.

\subsection{Filter-Bank Scale Effects}

Filter-bank scale (ascendant, constant, or descendant frequency weighting)
affects primarily the LFCC and MFCC front-ends.  Descendant filter banks,
which place greater weight on low frequencies, consistently yield the lowest
median EER, consistent with the finding that human voice formants~F1 and~F2
reside below 900~Hz and carry the primary anti-spoofing cues.  The ascendant
bank performs worst across all configurations and loss functions.

LFCC with a constant filter bank and LCNN-LSTM-sum/P2S-Grad achieves the
best median EER of $2.58\%$ on the English test set, with the descendant bank
achieving a competitive median of $2.51\%$ but with slightly higher
inter-run variance.  Table~\ref{table:scale-best-eer} summarises the variance
statistics for the best-performing LFCC architecture across the three scales.

\begin{table}[htbp]
\centering
\caption{EER (\%) distribution for LCNN-LSTM-sum/P2S-Grad with LFCC across
         filter-bank scales (English eval, 6 runs per condition).}
\label{table:scale-best-eer}
\setlength\extrarowheight{-2pt}
\begin{tabular}{lcccccc}
\hline
\textbf{Scale} & \textbf{Max} & \textbf{Q3} & \textbf{Median} & \textbf{Q1} & \textbf{Min} \\
\hline
Ascendant  & 3.246 & 2.861 & 2.427 & 2.367 & 2.298 \\
Constant   & 2.763 & 2.639 & 2.577 & 2.444 & 2.148 \\
Descendant & 3.699 & 2.904 & 2.510 & 2.217 & 2.118 \\
\hline
\end{tabular}
\end{table}

LFCC also exhibits lower inter-run variance than MFCC across all scale
conditions, making it the more repeatable choice under random initialisation.
For practitioners prioritising stability, the constant filter bank combined
with LFCC and LCNN-LSTM-sum/P2S-Grad offers the best reliability.

%% ----------------------------------------------------------------
\section{Conclusions}

This paper presented a comprehensive, multi-condition evaluation of speech
anti-spoofing countermeasures exposing the impact of language, accent, and
filter-bank design on system performance.  The main conclusions are:

\begin{enumerate}
  \item \textbf{Cepstral front-ends dominate.}  LFCC and MFCC consistently
        outperform CQCC and Spectrogram.  LFCC achieves the best English EER
        ($2.67\%$) and lowest variance; MFCC achieves the best Spanish EER
        under bilingual training ($0.09\%$).

  \item \textbf{Language mismatch is severe without adaptation.}  Models
        trained exclusively on English degrade to $>30\%$ EER on Spanish,
        confirming that single-language training is insufficient for
        multilingual deployment.

  \item \textbf{Bilingual training achieves near-parity on both languages.}
        Merging English and Spanish training data yields near-zero EER on
        Spanish while sacrificing fewer than one percentage point on English,
        demonstrating that a single model can generalise across languages.

  \item \textbf{Optimal architecture is language- and accent-specific.}
        The best front-end/back-end combination varies by language and by
        Spanish regional accent, motivating accent-adaptive or ensemble
        strategies.

  \item \textbf{Descendant filter-bank scaling improves cepstral front-ends.}
        Placing greater weight on lower frequencies aligns with human voice
        formant structure and consistently reduces EER for both MFCC and LFCC.
\end{enumerate}

Future work will investigate transformer-based back-ends, self-supervised
front-end representations (wav2vec~2.0~\cite{baevski2020wav2vec},
HuBERT~\cite{hsu2021hubert}), deepfake detection with the ASVspoof~2021
dataset~\cite{delgado2021asvspoof}, and the role of speaker gender and age
on accent-based anti-spoofing robustness.

%% ----------------------------------------------------------------
\bibliographystyle{splncs04}
\bibliography{bib/local,bib/writing}

\end{document}
